% concatenated from new_paper1-3.tex
\documentclass[12pt,a4paper]{amsart}
\usepackage{amsmath, amssymb}
\usepackage{hyperref}
\usepackage[english]{babel}
\theoremstyle{plain}

\usepackage{cite}

\newtheorem{Theorem}{Theorem}
\newtheorem{Lemma}{Lemma}
\newtheorem{Proposition}{Proposition}
\newtheorem{Corollary}{Corollary}

\newtheorem{Example}{Example}
\newtheorem{Remark}{Remark}
\newtheorem{Conjecture}{Conjecture}

\begin{document}
	
	\title{Derivations of Lie Algebras of Vector Fields in Infinitely Many Variables}
	
	
		\author
	{Oksana Bezushchak}
	\address{Oksana Bezushchak: Faculty of Mechanics and Mathematics, Taras Shevchenko National University of Kyiv, Volodymyrska, 60, Kyiv 01033, Ukraine}
	\email{bezushchak@knu.ua}
	
	\author{Iryna Kashuba}
	\address{Iryna Kashuba: Shenzhen International Center for Mathematics, Southern University of Science and Technology, 1088 Xueyuan Avenue, Nanshan District, Shenzhen, Guangdong, China}
	\email{kashuba@sustech.edu.cn}    
	
%	\date{May 31, 2025.}
	\keywords{Algebra of infinite matrices,  derivation, Lie algebra, locally matrix algebra, polynomial algebra in infinitely many variables}
	\subjclass[2020]{Primary 16W25, 17B66, Secondary 17A36, 17B40,  17B56}
	\maketitle
	
		\begin{abstract}
	Let $W_X(\mathbb{F})$ be the Lie algebra of all derivations of the polynomial algebra $\mathbb{F}[X]$ in infinitely many variables.	We describe  all derivations of $W_X(\mathbb{F})$  over a field of  characteristic zero and  prove that all  such derivations are inner. We also consider   the subalgebras $W_\infty(\mathbb{F})$  and $W_{\text{fin}}(\mathbb{F})$ of the algebra $W_X(\mathbb{F})$ and describe  all of their derivations. 	
	\end{abstract}
	
	
	\section{Introduction} 
	
	Let $\mathbb{F}$ be a field of zero characteristic, and let $\mathbb{F}[X_n]$, $n \geq 1$, be the polynomial algebra on the set of variables $X_n = \{x_1, \dots, x_n\}$.
	
	The Lie algebra:
	\[
	W_n(\mathbb{F}) = \text{Der } \mathbb{F}[X_n] = \left\{ \sum_{i=1}^{n} f_i(x_1, \dots, x_n) \frac{\partial}{\partial x_i} \ \Big| \ f_i \in \mathbb{F}[X_n] \right\}
	\]
	of all derivations of the algebra $\mathbb{F}[X_n]$ is an infinite-dimensional simple Lie algebra of Cartan type (see \cite{Rudakov1,Strade}).
	
	Consider the ascending chains: \[
	X_1 \subset X_2 \subset \dots, \quad \text{and} \quad 
	W_1(\mathbb{F}) \subset W_2(\mathbb{F}) \subset \dots .
	\]
	Let $$X = \bigcup_{n \geq 1} X_n ,\quad \text{and} \quad  W_\infty(\mathbb{F}) = \bigcup_{n \geq 1} W_n(\mathbb{F}).$$ 
	
	The Lie algebra $W_\infty(\mathbb{F})$ is a subalgebra of the Lie algebra $W_X(\mathbb{F})$ of all derivations of the polynomial algebra $\mathbb{F}[X]$ in infinitely many variables:
	\[
	W_X(\mathbb{F}) = \text{Der } \mathbb{F}[X] = \left\{ \sum_{i=1}^{\infty} f_i(x_1, \ldots ) \frac{\partial}{\partial x_i} \ \Big| \ f_i \in \mathbb{F}[X] \right\}.
	\] We will consider also the Lie algebra $W_{\text{fin}}(\mathbb{F})$ lying strictly between $W_\infty(\mathbb{F})$ and $W_X(\mathbb{F})$. The algebra $W_{\text{fin}}(\mathbb{F})$ consists of derivations
	\[
	\sum_{i=1}^{\infty} f_i(x_1, \ldots ) \frac{\partial}{\partial x_i}
	\]
	such that each variable $x_n \in X$ occurs in only finitely many polynomials $f_i$, $i \geq 1$. It is easy to see that $W_\infty(\mathbb{F})$ is an ideal of the algebra $W_{\text{fin}}(\mathbb{F})$.
	
	
	Floris Takens \cite{Takens} (se also \cite{Farnsteiner,Morimoto}) proved that every derivation of the Lie algebra  $W_n(\mathbb{F})$ over a field of characteristic $0$ is inner. Richard E. Block, Robert L. Wilson, Helmut Strade, and Alexander Premet~\cite{Strade} studied $W_n(\mathbb{F})$ in positive characteristic. For more recent papers concerning algebraic structure of $W_n(\mathbb{F})$, see 	\cite{Bezushchak_Petravchuk_Zelmanov,Klymenko_Lysenko_Petravchuk,Makedonskyi_Petravchuk}.
	

	
	
	
	In this paper, we describe derivations of the algebra $W_\infty(\mathbb{F}),$ $W_{\text{fin}}(\mathbb{F}),$ and $ W_X(\mathbb{F}) $ over a field $\mathbb{F}$  of zero characteristic.
	
	\begin{Theorem}\label{Theo1} Every derivation of the algebra $W_\infty(\mathbb{F})$ is a restriction of an adjoint operator $\text{ad}(a):x\to [x,a]$, where $a \in W_{\text{fin}}(\mathbb{F})$. \end{Theorem}
	
	\begin{Theorem}\label{Theo2}  All derivations of the algebra $W_{\text{fin}}(\mathbb{F})$ are inner. \end{Theorem}
	
	
	\begin{Theorem}\label{Theo3} All  derivations of algebra $W_X(\mathbb{F})$ are inner. \end{Theorem}
	
	
	
	From Theorem \ref{Theo1}, it follows that not all derivations of $W_{\infty}(\mathbb{F})$ are inner, but the ideal of inner derivations is dense in the algebra of all derivations of $W_{\infty}(\mathbb{F})$ in the Tychonoff topology. Dragomir Ž. Doković and Kaiming Zhao  \cite{Doković_Zhao} proved that the ideal of inner derivations is dense in $W_X(\mathbb{F})$ in the Tychonoff topology in a more general context.
	
	
	
	
	There is an analogy between polynomial algebras in infinitely many variables and algebras of infinite matrices. Let $\mathbb{N}=\{1,2,\ldots\}$ be the set of all positive integers. We consider infinite $\mathbb{N} \times \mathbb{N}$ matrices over the field $\mathbb{F}$ of arbitrary characteristic. The Lie algebra $\mathfrak{gl}_\infty(\mathbb{F})$ consists of $\mathbb{N} \times \mathbb{N}$ matrices having finitely many nonzero elements, and let $$\mathfrak{sl}_\infty(\mathbb{F}) = [\mathfrak{gl}_\infty(\mathbb{F}), \mathfrak{gl}_\infty(\mathbb{F})].$$ Clearly, $$\mathfrak{sl}_\infty(\mathbb{F}) = \bigcup_{n=2}^{\infty} \mathfrak{sl}_n(\mathbb{F}).$$ The Lie algebra $\mathfrak{gl}_{\mathbb{N}}(\mathbb{F})$ consists of $\mathbb{N} \times \mathbb{N}$ matrices having finitely many nonzero entries in each column. Finally, the intermediate algebra $$\mathfrak{gl}_{\text{fin}}(\mathbb{F}), \quad  \mathfrak{gl}_{\infty}(\mathbb{F}) \subseteq \mathfrak{gl}_{\text{fin}}(\mathbb{F}) \subset \mathfrak{gl}_{\mathbb{N}}(\mathbb{F}),$$ consists of $\mathbb{N} \times \mathbb{N}$ matrices having finitely many nonzero entries in each row and in each column.
	
	
	
	
	C.-H. Neeb \cite{Neeb} showed that an arbitrary derivation of the Lie algebra $\mathfrak{sl}_\infty(\mathbb{F})$ is of the type $\text{ad}(a)$, $a \in \mathfrak{gl}_{\text{fin}}(\mathbb{F})$, provided that $\text{char} \, \mathbb{F} = 0$. In~\cite{Bezushchak2,Bezushchak_Kashuba_Zelmanov}, we extended this result to algebras of infinite matrices over fields of positive characteristics. We also proved that all derivations of the algebras $\mathfrak{gl}_{\text{fin}}(\mathbb{F})$ and $\mathfrak{gl}_N(\mathbb{F})$ are inner.
	
	The approach in \cite{Bezushchak2} was as follows:
	(1) to lift derivations of Lie algebras of infinite matrices to derivations of associative algebras of infinite matrices, using the solution to Herstein's problems by K.~Beidar, M.~Bre\v sar, M.~Chebotar, and W.S.~Martindale 3nd  (see, \cite{Bei_Bre_Cheb_Mart_1,Bei_Bre_Cheb_Mart_2,Bei_Bre_Cheb_Mart_3});
	(2)~to use N.~Jacobson's description of derivations of dense algebras of linear transformations \cite{Jacobson}.
	
	This approach does not work for Lie algebras of derivations of polynomials. For this reason, we will follow a different path, one that is closer to the description of derivations of locally matrix algebras; see~\cite{Bezushchak1,BezOl_2}.
	
	\section{On locally finite-dimensional $\mathfrak{sl}_\infty(\mathbb{F})$-modules}  
	
Let $\mathbb{N}$ be the set of positive integers.	Consider the Lie algebra:
	\[
	\text{span} \left\{ x_i \frac{\partial}{\partial x_j}   , \ i\not= j; \quad x_i \frac{\partial}{\partial x_i} - x_j \frac{\partial}{\partial x_j}, \ i, j \in \mathbb{N} \right\}.
	\]
	This algebra is isomorphic to $\mathfrak{sl}_\infty(\mathbb{F})$. Hence, we can view $\mathfrak{sl}_\infty(\mathbb{F})$ as a subalgebra of $W_\infty(\mathbb{F})$.  If we consider the Lie algebra \[
	\text{span} \left\{ x_i \frac{\partial}{\partial x_j}   ,  \ i, j \in \mathbb{N} \right\}
	\] that is isomorphic to $\mathfrak{gl}_\infty(\mathbb{F})$, then  the subalgebra $\mathfrak{sl}_\infty(\mathbb{F})$  is an ideal of the algebra $\mathfrak{gl}_\infty(\mathbb{F})$.
	
	
	Let $$X^* = \bigcup_{n= 0}^{\infty} \underbrace{X \cdots X}_n$$ be the set of all monomials in $X$; and let
	\[
	\widetilde{\mathbb{F}[X]} = \left\{ \sum_{m \in X^{*}} \alpha_m m \ \Big| \ \alpha_m \in \mathbb{F} \right\}
	\]
	be the vector space of infinite sums that can be identified with the vector space of all maps $X^* \to \mathbb{F}$. This vector space is in fact an algebra since every monomial can be expressed as a product of monomials in finitely many ways.
	
	Consider the vector space:
	\begin{equation}\label{eq1}
		\widetilde{W} = \sum_{i=1}^{\infty} \widetilde{\mathbb{F}[X]} \frac{\partial}{\partial x_i}.
	\end{equation}
	Infinite sums \eqref{eq1} cannot always be identified with derivations of the algebra $\widetilde{\mathbb{F}[X]}$, and a commutator of sums \eqref{eq1}  may not make sense. However, $\widetilde{W}$ can be viewed as a $W_\infty(\mathbb{F})$-module and, consequently, also as $\mathfrak{sl}_\infty(\mathbb{F})$-module, since for any $h_i\in \mathbb{F}[X]$ and $f_j \in \widetilde{\mathbb{F}[X]}$, we have:
	\[
	\Big[\sum_{i=1}^{n} h_i \frac{\partial}{\partial x_i}, \sum_{j=1}^{\infty} f_j \frac{\partial}{\partial x_j}\Big] = \sum_{i=1}^{n}   \Big[h_i \frac{\partial}{\partial x_i }, \sum_{j=1}^{\infty} f_j \frac{\partial }{\partial x_j}\Big] , \quad \text{where}\] 
	\[   \Big[h_i \frac{\partial}{\partial x_i}, \sum_{j=1}^{\infty} f_j \frac{\partial}{\partial x_j}\Big]= \sum_{j=1}^{\infty} h_i \Big(\frac{\partial f_j}{\partial x_i}\Big) \frac{\partial}{\partial x_j} -  \sum_{j=1}^{\infty} f_j \Big(\frac{\partial h_i}{\partial x_j}\Big) \frac{\partial}{\partial x_i} \in  \widetilde{W} ,\] i.e., $[ W_\infty(\mathbb{F}),\widetilde{W}] \subseteq \widetilde{W}.$ 
	
	Recall that a module $M$ over a Lie algebra $L$ is called \textit{locally finite-dimensional} if given a finite subset $S \subset M$ and a finitely generated subalgebra $L_1 \subset L$, the subset $S$ generates a finite-dimensional module over $L_1$.
	
	
	\begin{Lemma}\label{Lemma1} Let $L$ be a finite-dimensional Lie algebra over a field $\mathbb{F}$. Let $M_1, \dots, M_s$ be finite-dimensional $L$-modules, $s \geq 1$. Let $\{V_i: i \in I\}$ be a collection of $L$-modules, for some set of indices $I$, where each $V_i$ is isomorphic to one of the modules $M_1, \dots, M_s$. Then the Cartesian sum
		\[
		V = \bigoplus_{i \in I} V_i
		\]
		is a locally finite-dimensional module over $L$. \end{Lemma}
	
	
	\begin{proof} Let $a \in L$. Consider the linear transformations
		\[
		a\big|_{M_i} \in \text{End}_{\mathbb{F}}(M_i), \quad a\big|_{M_i}: x \mapsto ax, \quad x \in M_i.
	 \quad \text{for} \quad 1 \leq i \leq s.	\] Let $f_i(t)$ be the characteristic polynomial of the linear transformation $a|_{M_i}$. Then
		\[
		f_i(a|_{M_i}) = 0.
		\]
		
		Now set \( f(t) = f_1(t)\cdots f_s(t) \). Then $f(a\big|_{V_i}) = 0$ for any $i \in I$, and therefore $f(a\big|_V) = 0$. By the Poincaré–Birkhoff–Witt Theorem (see~\cite{Jacobson_PBW}), it implies that the module $V$ is locally finite-dimensional. This completes the proof of the lemma. \end{proof}
	
	For a monomial $m \in X^*$ of length $\kappa$, we define the degree $$\deg\Big(m \frac{\partial}{\partial x_i}\Big) = \kappa - 1.$$ Let $\widetilde{W}_\kappa$ denote the subspace $\widetilde{W}_\kappa \subset \widetilde{W}$ of elements of degree $\kappa$, $-1 \leq \kappa < \infty$. Clearly, $\widetilde{W}_\kappa$ is a $\mathfrak{sl}_{\infty}(\mathbb{F})$-submodule.
	
	\begin{Lemma}\label{Lemma2} For any $n \geq 1$, $-1 \leq \kappa < \infty$, the $\mathfrak{sl}_n(\mathbb{F})$-module $\widetilde{W}_\kappa$ is locally finite-dimensional. \end{Lemma}
	
	
	\begin{proof} Let $\mathbb{F}[X_n]_i$ be the linear span of all monomials in $X_n$ of length $i$; and let $U_i$ be the linear span of all monomials in $X \setminus X_n$ of length $i$ for $i\geq 1$. Then:
		\[
		\widetilde{W}_k = \sum_{1\leq j\leq n, \ 0 \leq i \leq k+1} \mathbb{F}[X_n]_i \ U_{k-i+1} \frac{\partial}{\partial x_j} + \sum_{j>n, \ 0 \leq i \leq k+1} \mathbb{F}[X_n]_i \ U_{k-i+1} \frac{\partial}{\partial x_j}
		\] for $-1 \leq k < \infty .$ This is a Cartesian sum of finite-dimensional $\mathfrak{sl}_n(\mathbb{F})$-modules,  $$\text{each isomorphic to} \quad \sum_{j=1}^{n} \ \mathbb{F}[X_n]_i \ \frac{\partial}{\partial x_j} \quad \text{or to} \quad \mathbb{F}[X_n]_i, \quad 0 \leq i \leq k+1.$$ By Lemma \ref{Lemma1}, the $\mathfrak{sl}_n(\mathbb{F})$-module $\widetilde{W}_k$ is locally finite-dimensional. This completes the proof of the lemma. \end{proof}
	
	The following consequence follows from Lemma \ref{Lemma2}.
	
	\begin{Corollary}\label{cor1} For  $-1  \leq k < \infty$, the $\mathfrak{sl}_\infty(\mathbb{F})$-module $\widetilde{W}_\kappa$ is locally finite-dimensional. \end{Corollary}
	
	Consider an element $w \in \widetilde{W}$,
	\[
	w = \sum_{l=1}^{\infty} f_l \frac{\partial}{\partial x_l}, \quad \text{where} \quad f_l \in \widetilde{\mathbb{F}[X]}.
	\]
	Let $1 \leq i \not= j < \infty$. Then
	\begin{equation}\label{eq2}
		\Big[w, x_j \frac{\partial}{\partial x_i}\Big] = f_j \frac{\partial}{\partial x_i} - \sum_{l=1}^{\infty} x_j \left(\frac{\partial f_l}{\partial x_i} \right) \frac{\partial}{\partial x_l}.
	\end{equation}
	
	The $\frac{\partial}{\partial x_i}$-component of the commutator above is:
	\begin{equation}\label{eq3}
		f_j - x_j \frac{\partial f_i}{\partial x_i}.
	\end{equation}
	
	Furthermore, we have:
	\[
	\Big[w, x_i \frac{\partial}{\partial x_i}\Big] = f_i \frac{\partial}{\partial x_i} - \sum_{l=1}^{\infty} \left(\frac{\partial f_l}{\partial x_i} \right) x_i \frac{\partial}{\partial x_l}, \quad \text{and}
	\]
	\[
	\Big[w, x_j \frac{\partial}{\partial x_j}\Big] = f_j \frac{\partial}{\partial x_j} - \sum_{l=1}^{\infty} \left(\frac{\partial f_l}{\partial x_j} \right) x_j \frac{\partial}{\partial x_l}.
	\] Hence:
	\begin{equation}\label{eq4}
		\Big[w, x_i \frac{\partial}{\partial x_i} - x_j \frac{\partial}{\partial x_j}\Big] = f_i \frac{\partial}{\partial x_i} - f_j \frac{\partial}{\partial x_j} + \sum_{l=1}^{\infty} \left(\frac{\partial f_l}{\partial x_j} x_j - \frac{\partial f_l}{\partial x_i} x_i \right) \frac{\partial}{\partial x_l}.
	\end{equation}
	
	
	
	
	
	\begin{Lemma}\label{Lemma3}  Let $n\geq 3.$ Then $$\Big\{w \in \widetilde{W} \ \Big|\  \mathfrak{sl}_n(\mathbb{F}) w = (0) \Big\} = $$ $$\widetilde{\mathbb{F}[X \setminus X_n]} \ \Bigg(\sum_{i=1}^{n} x_i \frac{\partial}{\partial x_i}\Bigg) + \sum_{l=n+1}^{\infty} \widetilde{\mathbb{F}[X \setminus X_n]} \frac{\partial}{\partial x_l}.$$ \end{Lemma}
	
	
	\begin{proof} Let 	\[
		w = \sum_{l=1}^{\infty} f_l \frac{\partial}{\partial x_l}, \quad \text{where} \quad f_l \in \widetilde{\mathbb{F}[X]} \quad \text{and} \quad \mathfrak{sl}_n(\mathbb{F}) w = (0) .
		\] From \eqref{eq2}, it follows that for any $l$ and $i,$ $l \geq n+1,$ $1 \leq i \leq n, $ we have
		\[
		\frac{\partial f_l}{\partial x_i} = 0, \quad  \text{ so } \quad f_l \in \mathbb{F}[X \setminus X_n].
		\]
		Thus, it follows that:
		\begin{equation}\label{eq4+}
			\mathfrak{sl}_n(\mathbb{F}) \ \Big(\sum_{i=1}^{n} f_i \frac{\partial}{\partial x_i} \Big) = (0).\end{equation}
		
		Let $1\leq i, j, l \leq n$ be distinct integers. By \eqref{eq4}, we have $$ \frac{\partial f_l}{\partial x_j} x_j =\frac{\partial f_l}{\partial x_i} x_i .$$
		
		
		For a monomial $m$, let $\deg_{x_i}(m)$ be the number of occurrences of $x_i$ in $m$. We showed that for every monomial $m$ appearing in $f_l$, and for any $1 \leq i \neq j \leq n$ with  $i, j$ distinct from $l$, the following holds:
		\[
		\deg_{x_i}(m) = \deg_{x_j}(m).
		\]
		It means that $$m = x_l^p (x_i \cdots x_{l-1} x_{l+1}  \cdots x_n)^q m',$$ where $m'$ is a monomial in $X \setminus X_n$. From the equality \eqref{eq2}, it follows that $q = 0$, $m = x_l^p m'$.
		
		The equality \eqref{eq3} implies that $f_i = x_i a_i$, where $a_i \in \widetilde{\mathbb{F}[X \setminus X_n]}$ for $1 \leq i \leq n$, and $a_i = a_j$ for all $1 \leq i, j \leq n$. This completes the proof of the lemma. \end{proof}
	
	The following consequence follows from Lemma \ref{Lemma3}.
	
	\begin{Corollary}\label{cor2} We have: $$ \big\{ w\in \widetilde{W} \ \big| \ \mathfrak{sl}_{\infty}(\mathbb{F}) w =(0) \big\}= \mathbb{F}\Big(\sum_{i=1}^{\infty} x_i \frac{\partial}{\partial x_i}\Big).$$ \end{Corollary}
	
	
	
	\begin{Lemma}\label{Lemma4} For any $n \geq 2$ and $-1 \leq k < \infty$, the first cohomology group vanishes:
		\[
		H^1\big(\mathfrak{sl}_n(\mathbb{F}), \widetilde{W}_k\big) = 0.
		\] \end{Lemma}
	
	
	\begin{proof}  Let $d: \mathfrak{sl}_n(\mathbb{F}) \to \widetilde{W}_k$ be a derivation; and  let $V$ be the $\mathfrak{sl}_n(\mathbb{F})$-submodule generated by $d(\mathfrak{sl}_n(\mathbb{F}))$. Since $\mathfrak{sl}_n(\mathbb{F})$ is perfect, i.e.,  $\mathfrak{sl}_n(\mathbb{F}) = [\mathfrak{sl}_n(\mathbb{F}) , \mathfrak{sl}_n(\mathbb{F})]$, it follows that $d(\mathfrak{sl}_n(\mathbb{F})) \subseteq V$. Lemma \ref{Lemma2} implies that $\dim_{\mathbb{F}} V < \infty$. Hence, the derivation $d: \mathfrak{sl}_n(\mathbb{F}) \to V$ is inner; see~\cite{Jacobson_PBW}. This completes the proof of the lemma. \end{proof}
	
	\begin{Lemma}\label{Lemma5} For any $-1 \leq k < \infty$, the first cohomology group vanishes: $$H^1\big(\mathfrak{sl}_{\infty}(\mathbb{F}), \widetilde{W}_k\big) = (0).$$\end{Lemma}
	
	
	\begin{proof} For an element $w \in \widetilde{W}_k$, let $\hat{w}$ denote the derivation defined by $$\hat{w}: x \mapsto xw, \quad  x \in \mathfrak{sl}_\infty(\mathbb{F}).$$ By Lemma \ref{Lemma4}, for any $n \geq 2$, there exists an element $w_n \in \widetilde{W}_k$ such that:
		\[
		(d - \hat{w}_n)(\mathfrak{sl}_n(\mathbb{F})) = (0).
		\]
		
		For an element $w \in \widetilde{W}$, a monomial $m\in X^{*}$ and   an integer $i\geq 1$, let $\gamma_w(m \frac{\partial}{\partial x_i})$ be the coefficient of $m \frac{\partial}{\partial x_i}$ in $w$. We have:
		\[
		\mathfrak{sl}_n(\mathbb{F})\left( \sum_{i=1}^{n} x_i \frac{\partial}{\partial x_i} \right) = (0);
		\] see \eqref{eq4+}.
		Considering elements $$w_n - \gamma_{w_n} \Big(x_1 \frac{\partial}{\partial x_1}\Big) \left( \sum_{i=1}^{n} x_i \frac{\partial}{\partial x_i} \right)$$ instead of $w_n$, we will assume that $$\gamma_{w_{n}}\Big(x_1 \frac{\partial}{\partial x_1}\Big) = 0 \quad \text{for any} \quad n \geq 2.$$
		We claim that for every monomial $m \in X^*$ and every integer $i\geq 1$, the sequence \begin{equation}\label{eq4++} \gamma_{w_2}\Big(m \frac{\partial}{\partial x_i}\Big), \quad \gamma_{w_3}\Big(m \frac{\partial}{\partial x_{i}}\Big), \quad \dots \end{equation} stabilizes. Indeed, let $m \in X^*_n$ and suppose $i<n$. Then $$\mathfrak{sl}_n(\mathbb{F}) \left( w_n - w_{n+1} \right) = (0).$$ Therefore, by Lemma \ref{Lemma3}, we have:
		\[
		w_n - w_{n+1} \in \widetilde{\mathbb{F}[X \setminus X_n]} \Big(\sum_{j=1}^{n} x_j \frac{\partial}{\partial x_j} \Big)+ \sum_{l=n+1}^{\infty} \widetilde{\mathbb{F}[X \setminus X_n]} \frac{\partial}{\partial x_l}.
		\]
		The term $m \frac{\partial}{\partial x_i}$ cannot appear on the right-hand side unless $m = x_i$. Hence, assuming $m \neq x_i$, we obtain::
		\[
		\gamma_{w_n}\Big(m \frac{\partial}{\partial x_i}\Big) = \gamma_{w_{n+1}}\Big(m \frac{\partial}{\partial x_i}\Big).
		\]
		Again, by Lemma \ref{Lemma3}:
		\[
		\gamma_{w_n -w_{n+1}}\Big(x_i \frac{\partial}{\partial x_i}\Big) = \gamma_{w_n -w_{n+1}}\Big(x_1 \frac{\partial}{\partial x_1}\Big)= 0,
		\]
		and sequence \eqref{eq4++} stabilizes.
		
		Define
		\[
		\gamma\left(m \frac{\partial}{\partial x_i}\right) = \lim_{n \to \infty} 
		\gamma_{w_n}\left(m \frac{\partial}{\partial x_i}\right).
		\]
		Consider the element
		\[
		w = \sum_{m \in X^*,\ i \geq 1} \gamma\left(m \frac{\partial}{\partial x_i}\right) 
		m \frac{\partial}{\partial x_i}.
		\]
		From the construction above, it follows that
		\[
		\mathfrak{sl}_n(\mathbb{F})(w - w_n) =( 0).
		\]
		Thus, for all \( a \in \mathfrak{sl}_n(\mathbb{F}) \), we have
		\[
		[a, w] = [a, w_n] = d(a).
		\] This completes the proof of the claim. \end{proof}
	
	\begin{Lemma}\label{Lemma6} The first cohomology group vanishes: $$H^1\big(\mathfrak{sl}_{\infty}(\mathbb{F}), \widetilde{W}\big) = (0).$$ \end{Lemma}
	\begin{proof} Let $d:\mathfrak{sl}_{\infty}(\mathbb{F})\to \widetilde{W}$ be a derivation. For an arbitrary element $a\in \mathfrak{sl}_{\infty}(\mathbb{F}),$ the image $d(a)$ is an (infinite) sum of homogeneous components: $d(a)=d_{-1}(a)+d_0(a)+\cdots , $ where $d_i(a)\in \widetilde{W}_i.$ It is easy to  see that for every integer $i \geq -1$, the mapping $a\to d_i(a)$ is a derivation from $\mathfrak{sl}_{\infty}(\mathbb{F})$ to $\widetilde{W}_i$.  By Lemma \ref{Lemma5}, there exists an element $w_i \in \widetilde{W}_i$ such that $d_i(a) = [a, w_i]$ for an arbitrary element $a \in \mathfrak{sl}_\infty(\mathbb{F})$. Thus, $$w = \sum_{i=-1}^{\infty} w_i, \quad d = \hat{w}.$$ This completes the proof of the lemma. \end{proof}
	
	\section{The poof of Theorem \ref{Theo1}}
	
	The Lie algebras $W_{\text{fin}}(\mathbb{F})$ and $W_X(\mathbb{F})$ can be naturally identified with subspaces of $\widetilde{W}$. We denote these subspaces as $\widetilde{W}_{\text{fin}}, \widetilde{W}_X$ respectively. Recall that
	\[
	\widetilde{W}_X = \left\{ \sum_{i=1}^{\infty} f_i \frac{\partial}{\partial x_i} \ \middle| \ f_i \text{ are polynomials in } X \right\}.
	\]
	
	
	
	\begin{Lemma}\label{Lemma7} Let $w \in \widetilde{W}$, and $\mathfrak{sl}_\infty(\mathbb{F}) w \subseteq \widetilde{W}_X$. Then
		\[
		w = a \sum_{i=1}^{\infty} x_i \frac{\partial}{\partial x_i} + w',
		\]
		where $w' \in \widetilde{W}_X$, and for an arbitrary $i$, the partial derivative $\frac{\partial a}{\partial x_i}$ is a polynomial. \end{Lemma}
	
	\begin{proof} Let $$w = \sum_{i=1}^{\infty} f_i \frac{\partial}{\partial x_i}, \quad \text{where} \quad f_i \in \widetilde{\mathbb{F}[X]}.$$ From \eqref{eq2}, it immediately follows that for any $1\leq   i \neq j < \infty$, the partial derivative $\frac{\partial f_i}{\partial x_j}$ is a polynomial.
		
		
		
		The $\frac{\partial}{\partial x_i}$-component of the right-hand side of \eqref{eq4} is the following: 
		\[
		f_i \frac{\partial}{\partial x_i} + \left( \frac{\partial f_i}{\partial x_j} x_j - \frac{\partial f_i}{\partial x_i} x_i \right) \frac{\partial}{\partial x_i}.
		\]
		Hence, 
		\[
		f_i - \frac{\partial f_i}{\partial x_i} x_i 
		\] is also a polynomial.
		
		Let $$f_i = x_i \cdot h_i +\bar{f}_i, \quad \text{where} \quad  \bar{f}_i \in \mathbb{F}[X \setminus \{x_i\}], \quad h_i \in \mathbb{F}[X].$$ Thus, $$ f_i - \frac{\partial f_i}{\partial x_i} x_i =   x_i \cdot h_i + \bar{f}_i -\Big(h_i + x_i \frac{\partial h_i}{\partial x_i}\Big)x_i = \bar{f}_i  - x_i^2 \ \frac{\partial h_i}{\partial x_i}$$ is a polynomial. It implies that $\bar{f}_i$ is also a polynomial. Let $h_i = h_i' + h_i''$, where $h_i'$ is the sum of monomials involving $x_i$, and 
		\[
		\frac{\partial h_i''}{\partial x_i} = 0.
		\] It follows from the above that $h_i'$ is also a polynomial. Then, summarizing, we get that:
		$f_i = x_i \cdot g_i + p_i,$  where $ p_i $  is a polynomial and $\frac{\partial g_i}{\partial x_i} = 0.$
		
		Remark that
		\[
		w - \sum_{i=1}^\infty x_i \ g_i \ \frac{\partial}{\partial x_i} \in \widetilde{W}_X.
		\]
		Therefore, consider $ \sum_{i=1}^\infty x_i \ g_i \ \frac{\partial}{\partial x_i} $ instead of \( w \), that is, from now on 
		\[
		w := \sum_{i=1}^\infty x_i g_i \frac{\partial}{\partial x_i}.
		\] It follows from \eqref{eq3} that 
		\[
		x_j g_j - x_
		j \frac{\partial (x_i g_i)}{\partial x_i} = x_j (g_j - g_i)
		\]
		is a polynomial.
		
		Let \( a = h_1 \). Then for an arbitrary \( i \), we have $ g_i = a + (\text{a polynomial}).$ Now
		\[
		w \in a \sum_{i=1}^\infty x_i \frac{\partial}{\partial x_i} + \widetilde{W}_X.
		\] Since $$  \frac{\partial (a x_j)}{\partial  x_i}, \quad  i \neq j,$$ is a polynomial, it follows that 
		$\frac{\partial a}{\partial x_i}$ is a polynomial. This completes the proof of the lemma. 
	\end{proof}
	
	
	Consider the subalgebra 
	\[
	L := \sum_{i=1}^\infty \mathbb{F} \frac{\partial}{\partial x_i} + \mathfrak{gl}_\infty(\mathbb{F})
	\]
	of  the algebra $ W_\infty(\mathbb{F}).$
	
	\begin{Lemma}\label{Lemma8} The first cohomology group vanishes: $$H^1(L, \widetilde{W}_X) = (0).$$ \end{Lemma}
	
	\begin{proof} Let \( d : L \to \widetilde{W}_X \) be a derivation. Let's show that it is inner. By Lemma \ref{Lemma6}, there exists an element \( w \in \widetilde{W} \) such that
		\[
		d(x) = [w, x] \quad \text{for any } x \in \mathfrak{sl}_\infty(\mathbb{F}).
		\] By Lemma \ref{Lemma7}, $$ w = a \sum_{i=1}^\infty x_i \frac{\partial}{\partial x_i} + w',$$ where $ w' \in \widetilde{W}_X , $ and  $\partial a / {\partial x_i}$ is a polynomial for every  $i. $
		
		
		Take \( 1 \leq i \not= p < \infty \), $p\not= q.$ Then
		\[
		\left[ \frac{\partial}{\partial x_i}, x_p \frac{\partial}{\partial x_q} \right] = 0.
		\] Applying the derivation \( d \), we get:
		\[
		\left[ d\left( \frac{\partial}{\partial x_i}\right), \ x_p \frac{\partial}{\partial x_q} \right]
		+ \left[ \frac{\partial}{\partial x_i}, \Big[w, x_p \frac{\partial}{\partial x_q}\Big] \right]
		= \] \[\left[ d\left( \frac{\partial}{\partial x_i} \right) + \Big[\frac{\partial}{\partial x_i}, w\Big], \  x_p \frac{\partial}{\partial x_q} \right] = 0.
		\]
		
		Let
		\[
		d\left( \frac{\partial}{\partial x_i} \right) = \sum_{k=1}^\infty u_k \frac{\partial}{\partial x_k},\] where  $ u_k$  are polynomials.  We obtain: 
		\[
		\left[ \frac{\partial}{\partial x_i}, w \right] = \left[ \frac{\partial}{\partial x_i}, \ a\sum_{j=1}^\infty x_j \frac{\partial}{\partial x_j} +   w' \right]=
		\]
		\[ \left( \frac{\partial a}{\partial x_i}\right)  \sum_{j=1}^\infty x_j \frac{\partial}{\partial x_j} + a  \frac{\partial }{\partial x_i} + \left[ \frac{\partial}{\partial x_i}, w' \right].  \]
		
		
		Let
		\[
		\left[ \frac{\partial}{\partial x_i}, w' \right] = \sum_k v_k \frac{\partial}{\partial x_k}, \] where   $v_k$ are polynomials. The \( \frac{\partial}{\partial x_i} \)-component of the element \[   d\left(\frac{\partial}{\partial x_i}\right) +\left[  \frac{\partial}{\partial x_i}, w\right] \quad \text{is} \quad 
		\left(u_i + \left(\frac{\partial a}{\partial x_i}\right) x_i + a + v_i\right) \frac{\partial}{\partial x_i}.
		\]
		
		
		Suppose that the element $a$ is not a polynomial. Since \[  u_i + \left(\frac{\partial a}{\partial x_i} \right)x_i + v_i \] is a polynomial and every variable is involved in finitely many monomials in \( a \), it follows that there exists \( q \neq i \) such that
		\[
		\frac{\partial}{\partial x_q}\left(u_i + \left(\frac{\partial a}{\partial x_i}\right) x_i + a + v_i\right) \neq 0.
		\]
		
		Let $p \neq i,$ and $ p \neq q.$  Then the element
		\[
		\left[  d\left( \frac{\partial}{\partial x_i} \right) + \left[ \frac{\partial}{\partial x_i}, w \right] , \ x_p \frac{\partial}{\partial x_q}\right]
		\]
		has a nonzero \( \frac{\partial}{\partial x_i} \)-component. We got a contradiction. Thus, the element $a$ is a polynomial, and $w\in \widetilde{W}_X.$  
		
		
		Now, considering the derivation $d - \hat{w}$ instead of $d$, we assume that $d(\mathfrak{sl}_\infty(\mathbb{F})) = (0).$ Let
		\[
		d\left(\frac{\partial}{\partial x_i}\right) = \sum_{j=1}^{\infty} f_j \frac{\partial}{\partial x_j}, \] where each $ f_j $ is a polynomial in $X.$ Arguing as above, we see that for $p \neq i$ and $q \neq p$:
		\begin{equation}\label{eq5}
			\left[ \sum_{j=1}^{\infty} f_j  \frac{\partial}{\partial x_j}, \ x_p \frac{\partial}{\partial x_q} \right] = 
			f_p \frac{\partial}{\partial x_q} - x_p  \sum_{j=1}^{\infty} \ \frac{\partial f_j}{\partial x_q} \ \frac{\partial}{\partial x_j} = 0 .
		\end{equation}
		Fix $j \geq 1$. Choosing $q \neq j$, we get $\frac{\partial f_j}{\partial x_q} = 0.$  Hence,
		$f_j = f(x_j)$ is a polynomial in $x_j.$
		
		Now the right-hand side of \eqref{eq5} looks as
		\[
		f_p  \frac{\partial}{\partial x_q} - x_p  \frac{\partial f_p}{\partial x_q}\  \frac{\partial}{\partial x_p} = 0,
		\]
		which implies
		\[
		f_p(x_p) = x_p  \frac{\partial f_q(x_q)}{\partial x_q}.
		\] It implies that $f_q(x_q) = \alpha_q + \beta_q x_q$, for $\alpha_q, \beta_q \in \mathbb{F},$ and $f_p(x_p) =  \beta_q x_p.$  Hence, there exists $\beta \in \mathbb{F}$ such that for every $p\not= i, $  we have 
		$f_p = \beta x_p.$ Choosing $q = i$, we get $f_i = \alpha + \beta x_i$. Thus,
		\[
		d\left( \frac{\partial}{\partial x_i} \right) = \alpha_i  \frac{\partial}{\partial x_i} + \beta_i  \sum_{j=1}^{\infty} x_j  \frac{\partial}{\partial x_j}, \quad \text{where} \quad \alpha_i, \ \beta_i \in \mathbb{F}.
		\]
		
		Let $1 \leq i \neq j < \infty$. Applying the derivation $d$ to $[\frac{\partial}{\partial x_i}, \frac{\partial}{\partial x_j}] = 0$, we get:
		\[
		\left[ \alpha_i   \frac{\partial}{\partial x_i} + \beta_i  \sum_{k=1}^{\infty}  x_k  \frac{\partial}{\partial x_k}, \ \frac{\partial}{\partial x_j} \right] + \left[ \frac{\partial}{\partial x_i}, \ \alpha_j  \frac{\partial}{\partial x_j} + \beta_j  \sum_{k=1}^{\infty}  x_k  \frac{\partial}{\partial x_k} \right] = 
		\]
		\[ \beta_j  \frac{\partial}{\partial x_i} -\beta_i \frac{\partial}{\partial x_j} = 0,
		\]
		which implies $\beta_i = \beta_j = 0$.
		Hence, $$d\left(\frac{\partial}{\partial x_i}\right) = \alpha_i   \frac{\partial}{\partial x_i} \quad \text{for every} \quad  i \geq 1.$$
		
		
		Applying the derivation $d$ to
		\[
		\left[ \frac{\partial}{\partial x_i}, \ x_i  \frac{\partial}{\partial x_q} \right] = \frac{\partial}{\partial x_q}, \quad i \neq q,
		\] we get: $$ \alpha_i   \frac{\partial}{\partial x_q}= \alpha_q   \frac{\partial}{\partial x_q}.$$ Hence $ \alpha_i =\alpha$  for all $ i\geq 1. $ Now
		\[
		d\left( \frac{\partial}{\partial x_i} \right) = \alpha  \frac{\partial}{\partial x_i} = \left[ \frac{\partial}{\partial x_i}, \ \alpha \sum_{k=1}^{\infty}  x_k  \frac{\partial}{\partial x_k} \right].
		\]
		
		
		Replace \( d \) with $  d - \widehat{u} ,$ where $$ u= \alpha \sum_{k=1}^\infty x_k \frac{\partial}{\partial x_k}.$$
		Since
		\[
		\left[\mathfrak{sl}_\infty(\mathbb{F}), \ \sum_{k=1}^\infty x_k \frac{\partial}{\partial x_k}\right] = (0),
		\]
		we still have \( d(\mathfrak{sl}_\infty(\mathbb{F})) = (0) \),  and in addition to that,
		\[
		d\left( \frac{\partial}{\partial x_i} \right) = 0 \quad \text{for any } i \geq 1.
		\]
		
		We have:
		\[
		\mathfrak{gl}_\infty(\mathbb{F}) = \sum_{i=1}^\infty \mathbb{F} x_i \frac{\partial}{\partial x_i} + \mathfrak{sl}_\infty(\mathbb{F}).
		\] Let us prove that \[
		d\left( x_i \frac{\partial}{\partial x_i}\right) = 0.
		\] 
		
		For any \( p \geq 1 \), we compute
		\[
		\left[ \frac{\partial}{\partial x_p}, \ x_i \frac{\partial}{\partial x_i} \right] =
		\begin{cases}
			\frac{\partial}{\partial x_i}, & \text{if } p = i, \\
			0, & \text{if } p \neq i.
		\end{cases}
		\] In both cases,
		\[
		d\left( \left[ \frac{\partial}{\partial x_p}, \ x_i \frac{\partial}{\partial x_i} \right] \right) = 0.
		\]
		Since $$ d\left( \frac{\partial}{\partial x_p} \right) = 0,  \quad \text{then} \quad 
		\left[ \frac{\partial}{\partial x_p},\  d \left( x_i \frac{\partial}{\partial x_i} \right) \right] = 0.$$
		Assume
		\[
		d \left( x_i \frac{\partial}{\partial x_i} \right) = \sum_{k=1}^{\infty} a_k \frac{\partial}{\partial x_k}.
		\]
		For any \( k \geq 1 \), we have \( \frac{\partial}{\partial x_p} a_k = 0 \). Hence, \( a_k \in \mathbb{F} \). From the fact that 
		\[
		d \left( x_i \frac{\partial}{\partial x_i} - x_j \frac{\partial}{\partial x_j} \right) = 0,
		\]
		it follows that the element 
		\[
		\sum_{k = 1}^{\infty} a_k \frac{\partial}{\partial x_k}
		\]
		does not depend on the choice of the index  \( i \).
		
		Let \( q \ne i \). Then
		\[
		\left[ x_i \frac{\partial}{\partial x_i}, \ x_i \frac{\partial}{\partial x_q} \right] = x_i \frac{\partial}{\partial x_q}.
		\] Applying the derivation \( d \), we obtain:
		\[
		\left[ \sum_{k=1}^\infty a_k \frac{\partial}{\partial x_k}, \ x_i \frac{\partial}{\partial x_q} \right] = 0.
		\] Consequently, \( a_i = 0 \), which completes the proof that
		\[
		d(\mathfrak{gl}_\infty(\mathbb{F})) = (0).
		\]
		
		Thus, \[
		d(a) = \left[a, \ \alpha  \sum_{k=1}^\infty  x_k  \frac{\partial}{\partial x_k}\right] = 0 \quad \text{for all} \quad a \in \mathfrak{gl}_\infty(\mathbb{F}).
		\] The element $$\alpha  \sum_{k=1}^\infty x_k  \frac{\partial}{\partial x_k}$$ lies in $\widetilde{W}_{\text{fin}}.$ Thus, we proved that $d$ is an inner derivation. This completes the proof of the lemma. 
	\end{proof}
	
	
	\begin{Lemma}\label{Lemma9} If $w \in \widetilde{W}_X$ and $\mathfrak{sl}_\infty(\mathbb{F}) w \subseteq \widetilde{W}_{\text{fin}}$, then $w \in \widetilde{W}_{\text{fin}}$. \end{Lemma} 
	
	\begin{proof} Let $$w = \sum_{j=1}^{\infty} f_j \frac{\partial}{\partial x_j}\in \widetilde{W}_X.$$ Suppose $p \neq q$. By \eqref{eq5}, we get:
		\[
		w' = f_p \frac{\partial}{\partial x_q}- x_p \sum_{j=1}^{\infty} \left( \frac{\partial f_j}{\partial x_q}\right) \frac{\partial}{\partial x_j} \in \widetilde{W}_{\text{fin}}.
		\] If the variable $x_q$ is involved in $f_j$ for $j \neq q$, then the variable $x_p$ is involved in the $\frac{\partial}{\partial x_j}$-component of $w'$. Since $w' \in \widetilde{W}_{\text{fin}}$, it follows that $x_q$ is involved in finitely many $f_j$'s. This completes the proof of the lemma. \end{proof}
	
	It remains to prove the following lemma.
	\begin{Lemma}\label{Lemma10}  Let $d : W_\infty(\mathbb{F}) \to \widetilde{W}$ be a derivation such that
		\[
		d\left( \sum_{i=1}^{\infty} \mathbb{F}  \frac{\partial}{\partial x_i} + \mathfrak{gl}_\infty(\mathbb{F}) \right) = (0).
		\]
		Then \( d = 0 \). \end{Lemma}
	
	\begin{proof} Let \( m \) be a monomial. We will use induction on the length of the monomial \( m \) and show that
		\[
		d\left( m \frac{\partial}{\partial x_i} \right) = 0.
		\]  
		If the length is $0$ or $1,$ then everything clearly follows from the assumption. Suppose that  the length of \( m \) be greater than or equal to  \(  2 \).  
		Let
		\[
		d\left( m \frac{\partial}{\partial x_i} \right) = \sum_{j=1}^\infty f_j \frac{\partial}{\partial x_j}.
		\]  
		By the induction assumption, for any \( k \geq 1 \), we have:
		\[
		d\left( \Big[m \frac{\partial}{\partial x_i}, \  \frac{\partial}{\partial x_k}\Big] \right) 
		= \sum_{j=1}^\infty   \Big( \frac{\partial f_j}{\partial x_k} \Big) \frac{\partial}{\partial x_j}  = 0.
		\]  
		This implies that \( f_j \in \mathbb{F} \) for any \( j \geq 1 \).
		
		Now let \( p \neq i \). Choose \( q \geq 1 \) such that \( p \neq q \) and  \(  x_q \) does not appear  in the monomial \( m \). Then we have that 
		\[
		\left[m \frac{\partial}{\partial x_i}, \ x_p \frac{\partial}{\partial x_q}\right] = 0,
		\]
		 which implies the following equality:
		\begin{align*}
			d\left( \Big[m \frac{\partial}{\partial x_i}, x_p \frac{\partial}{\partial x_q}\Big] \right) 
			&= \Big[d\Big(m \frac{\partial}{\partial x_i}\Big),  x_p \frac{\partial}{\partial x_q}\Big] + \Big[m \frac{\partial}{\partial x_i},  d\Big(x_p \frac{\partial}{\partial x_q}\Big)\Big] = 0.
		\end{align*}
		
		Moreover, 
		\[ x_p \frac{\partial}{\partial x_q} \in \mathfrak{sl}_\infty(\mathbb{F}),
		\quad \text{so} \quad d\left(x_p \frac{\partial}{\partial x_q}\right) = 0.
		\]
		Consequently, 
		\[
		\left[d\left(m \frac{\partial}{\partial x_i}\right), \ x_p \frac{\partial}{\partial x_q}\right] = \left[ \sum_{j=1}^{\infty} f_j \frac{\partial}{\partial x_j}, \ x_p \frac{\partial}{\partial x_q} \right] =  f_p \frac{\partial}{\partial x_q}=0.
		\]
		It follows that \( f_q = 0 \) for all $p\not= i$; and  
		\[
		d\left( m \frac{\partial}{\partial x_i} \right) = \alpha \frac{\partial}{\partial x_i}, \quad \alpha \in \mathbb{F}.
		\]
		
		First, suppose that \( x_i \) is involved in \( m \); that is, \( m = x_i m' \).  
		Choose \( j \) such that \( x_j \) is not involved in \( m \). Then:
		\[
		\left[m' x_j \frac{\partial}{\partial x_i}, \ x_i \frac{\partial}{\partial x_j}\right] 
		= m' x_j \frac{\partial}{\partial x_j} - x_i m' \frac{\partial}{\partial x_i}.
		\]
		Applying the derivation \( d \) to both sides, we obtain:
		\[
		\left[d\Big(m' x_j \frac{\partial}{\partial x_i}\Big), \ x_i \frac{\partial}{\partial x_j}\right]
		= d\Big(m' x_j \frac{\partial}{\partial x_j}\Big) - d\Big( m \frac{\partial}{\partial x_i}\Big).
		\]
		The left-hand side lies in \( \mathbb{F}  \frac{\partial}{\partial x_j} \), 
		as does the right-hand side: $d\Big(m' x_j \frac{\partial}{\partial x_j}\Big) $ lies in $ \mathbb{F}  \frac{\partial}{\partial x_j} $, whereas 
		$d\Big( m \frac{\partial}{\partial x_i}\Big)$ lies in $\mathbb{F}  \frac{\partial}{\partial x_i}.$  Hence, $
		d\left(m \frac{\partial}{\partial x_i}\right) = 0.$
		
		Now suppose that \( x_i \) is not involved in \( m \).  We have: $ m = m' x_k$ for $k\not= i.$ 
		Then:
		\[
		\left[x_k \frac{\partial}{\partial x_i},\ m' x_i \frac{\partial}{\partial x_i}\right]
		= m \frac{\partial}{\partial x_i}.
		\]
		By the above,
		\[
		d\left(m' x_i \frac{\partial}{\partial x_i}\right)= 0, \quad \text{which implies that:} \quad 
		d\left(m \frac{\partial}{\partial x_i}\right) = 0.
		\]
		Thus,  \( d = 0 \). This completes the proof of the lemma. \end{proof}
	
	
	
	\begin{proof}[Proof of Theorem \ref{Theo1}]
		Let \( d: W_\infty(\mathbb{F}) \to W_\infty(\mathbb{F}) \) be a derivation.  
		By Lemmas \ref{Lemma8} and \ref{Lemma9}, there exists an element \( w \in W_{\mathrm{fin}}(\mathbb{F}) \) such that 
		$d(a) = [a, w]$  for any element $$ a \in \sum_{i=1}^\infty \mathbb{F}  \frac{\partial}{\partial x_i} + \mathfrak{gl}_\infty(\mathbb{F}).$$ Now by Lemma \ref{Lemma10}, we conclude that \( d = \hat{w} \).  
		This completes the proof of Theorem \ref{Theo1}. \end{proof}
	
	
	
	\section{The poof of Theorem \ref{Theo2}}
	
	\begin{Lemma}\label{Lemma11} The centralizer of \( W_\infty(\mathbb{F}) \) in \( \widetilde{W} \) is zero. \end{Lemma}
	
	\begin{proof} Suppose that \( w \in \widetilde{W} \) and $[W_\infty(\mathbb{F}), w] = (0).$ 
		In particular,
		\[
		\left[ \frac{\partial}{\partial x_i},\ w\right] = 0 \quad \text{for every } i \geq 1.
		\]
		This implies that 
		\[
		w = \sum_{i=1}^\infty \alpha_i \frac{\partial}{\partial x_i}, \quad \text{where} \quad \alpha_i \in \mathbb{F}.
		\]
		Let us assume that \( \alpha_i \neq 0 \). Then
		\[
		[w, x_i \frac{\partial}{\partial x_1}] = \alpha_i \frac{\partial}{\partial x_1}  \neq 0.
		\]
		This contradicts the assumption and completes the proof of the lemma. \end{proof}
	
	
	
	\begin{proof}[Proof of Theorem \ref{Theo2}] Let \( d : W_{\mathrm{fin}}(\mathbb{F}) \to W_{\mathrm{fin}}(\mathbb{F}) \) be a derivation. Arguing as in the proof of Theorem \ref{Theo1}, we find an element \( w \in W_{\text{fin}}(\mathbb{F}) \) such that  
		\[
		(d - \hat{w})(W_{\infty}(\mathbb{F})) = (0).
		\]
		Since \( W_{\infty}(\mathbb{F}) \) is an ideal in \( W_{\text{fin}}(\mathbb{F}) \), we conclude that
		\[
		(d - \hat{w})\Big([W_{\text{fin}}(\mathbb{F}), W_{\infty}(\mathbb{F})]\Big) = \Big[(d - \hat{w}) \big( W_{\text{fin}}(\mathbb{F})\big), \ W_{\infty}(\mathbb{F}) \Big]= (0).
		\]
		By Lemma \ref{Lemma11},  we have that \( d = \hat{w} \), which completes the proof of  Theorem \ref{Theo2}. \end{proof}
	
	
	
	\section{The poof of Theorem \ref{Theo3}}
	
	To prove Theorem \ref{Theo3}, we will need a couple of lemmas. 
	
	Arguing as above, we assume that $d:W_X (\mathbb{F}) \to W_X (\mathbb{F})$ is a derivation and $d(W_{\infty}(\mathbb{F}))=(0).$
	\begin{Lemma}\label{Lemma12} Let \( i_0 > 1 \), 
		\[ \text{and let} \quad 
		a \in \sum_{j \neq i_0} \mathbb{F}[X \setminus \{x_{i_0}\}] \frac{\partial}{\partial x_j}. \quad \text{Then} \quad   d(a) \in \sum_{j \neq i_0} \mathbb{F}[X \setminus \{x_{i_0}\}] \frac{\partial}{\partial x_j} .\]
	\end{Lemma}
	
	\begin{proof} Let 
		\[
		d(a) = \sum_{j=1}^{\infty} f_j \frac{\partial}{\partial x_j}.
		\]
		From \( [a, \frac{\partial}{\partial x_{i_0}}] = 0 \), it follows that
		\[
		\left[d(a), \ \frac{\partial}{\partial x_{i_0}}\right] = 0. \] Hence, $ f_j \in \mathbb{F}\left[X \setminus \{x_{i_0}\}\right]$ for all $ j \geq 1.$
		Similarly, 
		\[
		\left[d(a), \ x_{i_0} \frac{\partial}{\partial x_{i_0}}\right] = 0 .\] Thus, $f_{i_0} = 0,$ and the lemma is proved.  \end{proof}
	
	
	\begin{Lemma}\label{Lemma13} Let \[ a = \sum_{j \neq i_0} f_j \frac{\partial}{\partial x_j} \in W_X(\mathbb{F}), \] where all polynomials  \( f_j \) are divisible by \( x_{i_0}^k \) for some \( k \geq 1 \). Then 
		\[
		d(a) \in x_{i_0}^{k-1} \ W_X(\mathbb{F}).
		\] \end{Lemma}
	
	\begin{proof} Let \( f_j = x_{i_0}^k f'_j \), where \( f_j' \in \mathbb{F}[X] .\) Choose polynomials \( \widetilde{f_j} \in \mathbb{F}[X] \)  such that 
		$$\frac{\partial}{\partial x_{i_0}} \widetilde{f_j} = f_j' .$$ Then 
		\[
		a = \left[x_{i_0}^k \frac{\partial}{\partial x_{i_0}}  , \ \sum_{j \neq i_0} \widetilde{f_j} \frac{\partial}{\partial x_j}\right].
		\]
		This implies that
		\[
		d(a) \in  \left[x_{i_0}^k \frac{\partial}{\partial x_{i_0}}  , \ W_X(\mathbb{F})\right] \subseteq x_{i_0}^{k-1} \ W_X(\mathbb{F}),
		\]
		and completes the proof of the lemma. \end{proof}
	
	
	
	\begin{proof}[Proof of Theorem \ref{Theo3}] Suppose that $w \in W_X(\mathbb{F})$, and  \[ d(w) = \sum_{i=1}^{\infty} h_i \frac{\partial}{\partial x_i}  \neq 0 .\] Choose \( i_0 \geq 1 \) such that \( h_{i_0} \neq 0 \). Let \( l \geq 1 \) be such that \( h_{i_0} \notin x_{i_o}^l \mathbb{F}[X] \).
		
		Let 
		\[
		w = \sum_{i=0}^\infty f_i \frac{\partial}{\partial x_i}.
		\]
		Since \( d(W_{\infty}(\mathbb{F})) = (0) ,\) we may assume without loss of generality that \( f_{i_0} = 0 \). Then the element \( w \) can be represented as
		\[
		w = w_0 + x_{i_0} w_1 +  \cdots + x_{i_0}^l w_l + w', \quad \text{where}
		\]
		\[  w_0, \ w_1, \ \dots, \ w_l \in \sum_{j \neq i_0} \mathbb{F}[X \setminus \{x_{i_0}\}]  \frac{\partial}{\partial x_j}, \quad \text{and} \quad w' \in x_{i_0}^{l+1} \ W_X(\mathbb{F}) .\]
		
		Choose  $i_1 \neq i_0.$ Now let 
		\[
		w_k = \sum_{j \neq i_0} f_{k j}  \frac{\partial}{\partial x_j} \quad \text{for} \quad 0 \leq k \leq l.
		\] Choose polynomials $$\widetilde{f_{kj}} \in \mathbb{F}[X \setminus \{x_{i_0}\}] \quad \text{such that} \quad  \frac{\partial \widetilde{f_{kj}}}{\partial x_{i_1}}=f_{kj}.$$ Consider $$ 
		\widetilde{w_k} = \sum_{j \neq i_0} \widetilde{f_{k j}}  \frac{\partial}{\partial x_j} .$$ We get
		\[
		\left[ x_{i_0}^k \frac{\partial}{\partial x_{i_1}}, \ \widetilde{w_k} \right] = x_{i_0}^k w_k,\] which implies that \[  d( x_{i_0}^k\, w_k )= \left[d\left( x_{i_0}^k \frac{\partial}{\partial x_{i_1}} \right),  \widetilde{w_k} \right] +\left[ x_{i_0}^k \frac{\partial}{\partial x_{i_1}},  d(\widetilde{w_k}) \right]=  \left[ x_{i_0}^k \frac{\partial}{\partial x_{i_1}},  d(\widetilde{w_k}) \right].
		\]
		Hence, the \( \frac{\partial}{\partial x_{i_0}} \)-component of \( d\left( x_{i_0}^k w_k\right) \) is zero by Lemma \ref{Lemma12}. Since $d(w)=d(w_0)+d(x_{i_0}w_1)+\cdots + d(x_{i_0}^l w_l)+ d(w'),$ we get that the \( \frac{\partial}{\partial x_{i_0}} \)-component of \( d(w) \) is equal to the \( \frac{\partial}{\partial x_{i_0}} \)-component of \( d(w') \). By Lemma~\ref{Lemma13}, this $  \frac{\partial}{\partial x_{i_0}}$-component lies in $$ x_{i_0}^l\mathbb{F}[X] \frac{\partial}{\partial x_{i_0}} ,$$ a contradiction. Thus, we have shown that \( d = 0 \), which completes the proof of Theorem \ref{Theo3}. \end{proof}
	
	\section{Some Conjectures}
	
	
	
	
	A.~Rudakov \cite{Rudakov1} (see also V.~Bavula \cite{Bavula,Bavula2,Bavula3}, H.-P.~Kraft and A.~Re\-ge\-ta~\cite{Kraft}) showed that every automorphism of the Lie algebra $W_{n}(\mathbb{F})$ is a conjugation by an automorphism of $\mathbb{F}[X_n]$. Consider the following group of automorphisms of the polynomial algebra $\mathbb{F}[X]$: \\
	$\mathrm{Aut}_{\mathrm{fin}}\, \mathbb{F}[X] = \big\{ \varphi \in \mathrm{Aut}\, \mathbb{F}[X] \  \big| $   every variable $x_n $     occurs in finitely many polynomials $\varphi(x_i)$ for $i \geq 1,$   and in finitely many polynomials 	$\varphi^{-1}(x_i)$ for $ i \geq 1 \}.$

	
\begin{Conjecture} Every automorphism of the Lie algebra $W_{\infty}(\mathbb{F})$ is a conjugation by an automorphism $\varphi \in \mathrm{Aut}_{\mathrm{fin}}\, \mathbb{F}[X]$. \end{Conjecture}
	

\begin{Conjecture}  Every automorphism of the Lie algebra $W_{X}(\mathbb{F})$ is a conjugation by an automorphism from $\mathrm{Aut}\, \mathbb{F}[X]$. \end{Conjecture}

	
	
	
	\begin{thebibliography}{99}
		
		\bibitem{Bavula} Bavula V.V., The group of automorphisms of the Lie algebra of derivations of a polynomial algebra, J. Algebra Appl. \textbf{16}(05), 1750088 (2017).
		
		\bibitem{Bavula2} Bavula V.V., The group of automorphisms of the Lie algebra of derivations of a field of rational functions, Glasg. Math. J. \textbf{59}(3) (2017), P.513-524.
		
		\bibitem{Bavula3} Bavula V.V., The groups of automorphisms of the Lie algebras of triangular  polynomial derivations, J. Pure Appl. Algebra \textbf{218}(5) (2014), P.829-851.
		
		\bibitem{Bei_Bre_Cheb_Mart_1} Beidar K.I., Bre\v{s}ar M., Chebotar M.A.,  Martindale 3rd W.S.,  On Herstein's Lie map conjectures I, Trans. Amer. Math. Soc. \textbf{353} (2001), P.4235-4260.
		
		\bibitem{Bei_Bre_Cheb_Mart_2} Beidar K.I., Bre\v{s}ar M., Chebotar M.A.,  Martindale 3rd W.S.,  On Herstein's Lie map conjectures II, J. Algebra \textbf{238} (2001), P.239-264.
		
		\bibitem{Bei_Bre_Cheb_Mart_3} Beidar K.I., Bre\v{s}ar M., Chebotar M.A.,  Martindale 3rd W.S.,  On Herstein's Lie map conjectures III, J. Algebra \textbf{249} (2002), P.59-94.
		
		
		
		\bibitem{Bezushchak1} Bezushchak O., Derivations and automorphisms of locally matrix algebras, J.~Algebra  \textbf{576} (2021), P.1-26. 
		
		
		
		
		\bibitem{Bezushchak2} Bezushchak O., Automorphism and derivations of algebras of infinite matrix. Linear Algebra App. \textbf{650}(2) (2022), P.42-59.
		

		
		\bibitem{Bezushchak_Kashuba_Zelmanov} Bezushchak O., Kashuba I., Zelmanov E., On Lie isomorphisms of rings, Mediterr. J. Math. (2025) 22:80.
		
		\bibitem{BezOl_2} Bezushchak O., Oliynyk B., Primary decompositions of unital locally matrix algebras, Bull. Math. Sci. \textbf{10}(1) (2020), P.2050006-1 - 2050006-7.
		
		
		\bibitem{Bezushchak_Petravchuk_Zelmanov} Bezushchak O., Petravchuk A., Zelmanov E., Automorphisms and derivations of affine commutative and $PI$-algebras, Trans. Amer. Math. Soc. \textbf{377}(2) (2024), P.1335-1356 (2024).  
		
		\bibitem{Doković_Zhao}   Doković D.Ž.,  Zhao K., Derivations, Isomorphisms, and Second Cohomology of Generalized Witt Algebras, Trans. Amer. Math. Soc. \textbf{350}(2) (1998), P.643-664.
		
	\bibitem{Farnsteiner} Farnsteiner R., On the cohomology of associative algebras and Lie algebras,
	Proc. Amer. Math. Soc.  \textbf{99} (1987), P.415-420.
		
		\bibitem{Jacobson} Jacobson~N., Lectures in Abstract Algebra: \textbf{II}. Linear Algebra (Graduate Texts in Mathematics), Springer-Verlag, Berlin-Heidelberg-New York, 1975. %chapter IX.5, theorem 2, p.247
		
		\bibitem{Jacobson_PBW} Jacobson N., Lie Algebras. Interscience, New York, 1962.
		
		
		
		\bibitem{Klymenko_Lysenko_Petravchuk}  Klymenko I.S.,  Lysenko S.V.,  Petravchuk A.P., Lie algebras of derivations with large abelian ideals, Algebra Discrete Math. \textbf{28}(1) (2019), P.123-129.
		
	
		

		
				
		
	\bibitem{Kraft} 	Kraft H., Regeta A., Automorphisms of the Lie algebra of vector fields, J. Eur. Math. Soc. \textbf{19} (2017), P.1577-1588.
		
		\bibitem{Makedonskyi_Petravchuk} Makedonskyi Ie.O., Petravchuk A.P., On nilpotent and solvable Lie algebras of derivations, J. Algebra, \textbf{401} (2014), P.245-257.
		
			\bibitem{Morimoto} Morimoto T., The derivation algebras of the classical infinite Lie algebras, Math.	Kyoto Univ. \textbf{16} (1976), P.1-24.
		
		\bibitem{Neeb} Neeb K.-H., Derivations of locally simple Lie algebras, J. Lie Theory \textbf{15} (2005), P.589-594.
		
		

		
		
		\bibitem{Rudakov1}  Rudakov A.N., Groups of automorphism of infinite-dimensional simple Lie algebras, Mathematics of the USSR-Izvestiya \textbf{3}(4) (1969), 748-764.    
		
		
		


		\bibitem{Strade} Strade H., Simple Lie Algebras over Fields of Positive Characteristic, I. Structure Theory, de Gruyter Expositions in Mathematics \textbf{38}, Berlin-New York, 2004.
		
\bibitem{Takens}	Takens F., Derivations of vector fields, Compositio Mathematica \textbf{26} (1973), P.151-158.


	
	\end{thebibliography}
	
	
\end{document}


16W25 Derivations, actions of Lie algebras
17A36 Automorphisms, derivations, other operators
17B40 Automorphisms, derivations, other operators
17B56 Cohomology of Lie (super)algebras

17B66 Lie algebras of vector fields and related (super)
algebras



\begin{enumerate}
	\item   The question remains whether our algebras $W_{\text{fin}}(\mathbb{F}) $  and $ W_X(\mathbb{F}) $ are simple Lie algebras. 
	
	
	Similar question from \cite{Doković_Zhao} If $F$ is a field of characteristic $0$ and $A$ a torsion-free group of infinite
	rank, is it true that $Der(FA)$  is a simple Lie algebra?
	
	
	
	\item  J\"org Feldvoss, Salvatore Siciliano, Lie algebras whose derivation algebras are simple, arXiv:2501.15560v1 [math.RA] 26 Jan 2025 
	
	МОЖЛИВО посилання на цю статтю добавити
	
	It is well known that a finite-dimensional Lie algebra over a field of characteristic zero is simple exactly when its derivation algebra is simple. In this paper, we characterize those Lie algebras of arbitrary dimension over any field that have a simple derivation algebra. As an application we classify the Lie algebras that have a complete simple derivation algebra and are either finite-dimensional over an algebraically closed field of prime characteristic 
	$p>3$  or $Z$-graded of finite growth over an algebraically closed field of characteristic zero.
	
	
\end{enumerate}












Можливо ще прокоментувати:
\begin{enumerate}
	\item У диференціювання \( d: \mathfrak{sl}_\infty(\mathbb{F}) \rightarrow \widetilde{W} \)
	існує \( w \in \widetilde{W} \), таке що \( d(x) = [x, w] \);
	\item Якщо \( a \in \widetilde{W} \), то \( [\mathfrak{sl}_\infty(\mathbb{F}), a] \subseteq W_\infty(\mathbb{F}) \), 
	отже \( a \in W_{\text{fin}}(\mathbb{F}) \).
\end{enumerate}
